%-----------------------------------------------------------------------------
% Abstrak (Bahasa Indonesia)
%-----------------------------------------------------------------------------


\chapter*{Abstrak}

	Penelitian ini mengembangkan sistem berbasis \textit{cloud} berbiaya rendah untuk mendeteksi kesiapan panen padi, dirancang bagi petani gurem dan divalidasi pada skala \textit{bench} dengan empat tanaman \textit{polybag} di dalam satu \textit{litter box}. Sistem menyasar dua kesenjangan: ketidaksesuaian biaya pencitraan satelit dan wahana udara tanpa awak pada skala sub-hektare, serta keterbatasan \textit{closed-world} pengklasifikasi terlatih yang tetap memberi label fenologi pada citra non-padi atau yang tidak layak ditafsirkan. Modul ESP32-CAM AI-Thinker dengan sensor OV3660 menangkap satu citra kanopi dari atas pada tiap bangun \textit{deep-sleep} dan mengunggahnya melalui Google Cloud Storage ke layanan Flask di Google Cloud Run, yang menuliskan prediksi ke Firebase \textit{Realtime Database} untuk dasbor \textit{Progressive Web App} yang ditempatkan di Netlify. Empat keluarga klasifikasi dibandingkan pada dataset terkurasi 1140 citra (1036 untuk validasi silang \textit{stratified 5-fold}, 104 sisanya \textit{held-out test}) yang mencakup fase vegetatif, generatif, dan matur: \textit{machine learning} klasik atas indeks vegetasi warna dengan fitur tekstur \textit{Gray Level Co-occurrence Matrix}, \textit{backbone} konvolusional \textit{transfer learning}, \textit{hybrid late-fusion}, dan \textit{endpoint vision-language} yang dipakai sebagai penjaga gerbang \textit{open-world}, bukan pengklasifikasi utama. Konfigurasi \textit{hybrid} EfficientNetB0 dengan fitur ExGR memimpin Macro F1 uji pada 0.9662 dan \textit{Support Vector Machine} klasik dengan ExGR menyusul di 0.9577 dengan artefak 55 KB, tetapi kesetaraan ini tidak terbawa ke sensor ESP32-CAM yang di-\textit{deploy}, sehingga \textit{pipeline} tetap memakai pengklasifikasi \textit{hybrid}. Penjaga gerbang \textit{vision-language} yang di-\textit{deploy} (Gemini 2.5 Flash Lite) menolak seluruh dua belas citra non-padi pada \textit{false-accept rate} 0 persen dan \textit{false-reject rate} 4.81 persen, dengan latensi \textit{end-to-end} bermedian 8.57 detik, jauh di bawah jadwal pengambilan citra tiap sepuluh menit. Tata kelola data ditanam langsung ke dalam arsitektur melalui residensi penyimpanan di wilayah Indonesia, pemisahan identitas perangkat dan petani, kontrol akses \textit{Identity and Access Management} \textit{least-privilege}, serta jejak audit yang selaras dengan Undang-Undang Nomor 19 Tahun 2016 tentang Informasi dan Transaksi Elektronik, Peraturan Pemerintah Nomor 71 Tahun 2019, Undang-Undang Nomor 27 Tahun 2022 tentang Pelindungan Data Pribadi, ISO/IEC 30141:2024, dan Rekomendasi ITU-T Y.2060.\\
\\
\textbf{Kata kunci:} fenologi padi, pencitraan proksimal, ESP32-CAM, \textit{vision-language model}, \textit{open-set recognition}, tata kelola data IoT, pertanian rakyat.

%-----------------------------------------------------------------------------
